% ===========================================================================
% Section: Experiments (placeholder for future results)
% ===========================================================================
\section{Experimental Setup}
\label{sec:experiments}

\subsection{Implementation Details}

\winlic{} is implemented in C++17 with CUDA 11.7 and built using CMake 3.18+.
The framework depends on LibTorch 2.0 (PyTorch C++ API), OpenCV 4.7 (with
CUDA and opencv\_contrib), TensorRT 8.6 (for SPNet depth completion),
yaml-cpp 0.6, and Eigen 3.4. The build supports both MSVC 2022 (Windows)
and GCC 9+ (Linux) with no platform-specific preprocessor directives
beyond CUDA architecture selection.

The SPNet depth completion engine is pre-compiled as a TensorRT plan file
(\texttt{.engine}) for two resolutions: 480$\times$640 and 512$\times$640.
The LPIPS model is exported as a TorchScript module (\texttt{lpips\_alex.pt}).

\subsection{Datasets}

The framework is designed to process data from the following sources:

\begin{itemize}
    \item \textbf{FAST-LIVO2 datasets}: Pre-recorded LiDAR-visual-inertial
    sequences with poses from the FAST-LIVO2 odometry front-end. The
    trajectory is exported in TUM format for offline processing.

    \item \textbf{MCD (Multi-Camera Dataset)}: Outdoor urban sequences
    with LiDAR, camera, and IMU data.

    \item \textbf{R3LIVE datasets}: RGB-LiDAR fusion sequences for
    indoor and outdoor scenes.

    \item \textbf{Custom Orbbec captures}: Live data from Orbbec Femto
    Mega cameras (LiDAR + iToF + stereo + IMU), recorded via the
    \texttt{OrbbecSource} and exported for offline processing.
\end{itemize}

\subsection{Metrics}

Reconstruction quality is evaluated using standard metrics:

\begin{itemize}
    \item \textbf{PSNR} (Peak Signal-to-Noise Ratio): measures pixel-level
    fidelity of rendered images.
    \item \textbf{SSIM} (Structural Similarity Index): measures perceptual
    quality via luminance, contrast, and structure.
    \item \textbf{LPIPS} (Learned Perceptual Image Patch Similarity):
    measures perceptual distance using deep features from AlexNet.
\end{itemize}

Metrics are computed on both training viewpoints (keyframes used during
optimization) and novel viewpoints (intermediate frames not used for
training), providing a measure of generalization.

\subsection{Baseline Comparison}

The primary baseline is the original \gaussianlic{} system running on
Linux with ROS. Since \winlic{} preserves the Gaussian Splatting algorithm
verbatim, we expect identical reconstruction quality (PSNR, SSIM, LPIPS)
when processing the same data. Differences may arise from:

\begin{itemize}
    \item Floating-point determinism between MSVC and GCC.
    \item Depth completion quality when using fused depth (LiDAR + iToF + stereo)
    versus LiDAR-only depth.
    \item Frame alignment differences between ROS timestamp-based and
    file-index-based synchronization.
\end{itemize}

\subsection{Runtime Analysis}

The offline processing mode eliminates the real-time constraint, allowing
the system to process each frame at full speed without buffer overflow
concerns. Expected runtime on an NVIDIA RTX~4090:
\begin{itemize}
    \item Depth fusion: $\sim$5\,ms/frame
    \item SPNet completion: $\sim$10\,ms/frame
    \item Gaussian extension: $\sim$50\,ms/keyframe
    \item Gaussian optimization: $\sim$500\,ms/keyframe (100 iterations)
\end{itemize}

\subsection{Ablation Studies}

Planned ablation studies include:
\begin{itemize}
    \item Effect of depth fusion: LiDAR-only vs.\ LiDAR+iToF vs.\ LiDAR+iToF+stereo.
    \item Effect of keyframe interval: $k \in \{1, 3, 5, 10\}$.
    \item Effect of Gaussian memory cap on quality vs.\ memory usage.
    \item Effect of SH degree on rendering quality.
\end{itemize}

\begin{table}[h]
\centering
\caption{Default hyperparameters (preserved from \gaussianlic).}
\label{tab:hyperparams}
\begin{tabular}{lc}
\toprule
Parameter & Value \\
\midrule
Image resolution & $640 \times 512$ \\
SH degree & 3 \\
Keyframe interval $k$ & 5 \\
Position learning rate & $1.6 \times 10^{-4}$ \\
Feature learning rate & $5.0 \times 10^{-3}$ \\
Opacity learning rate & $5.0 \times 10^{-2}$ \\
Scaling learning rate & $5.0 \times 10^{-3}$ \\
Rotation learning rate & $1.0 \times 10^{-3}$ \\
$\lambda_{\text{dssim}}$ & 0.2 \\
$\lambda_{\text{depth}}$ & 0.005 \\
Optimization iterations & 100/keyframe \\
Max depth & 20\,m \\
\bottomrule
\end{tabular}
\end{table}
